\thispagestyle{plain}
\setlength{\topmargin}{0.25in}
\begin{center}
    \Large \textbf{\uppercase{Abstract}}
\end{center}

\vspace{3\baselineskip}

\noindent
We study the problem of anomaly detection within the framework of universal outlier hypothesis testing. The goal is to identify sequences 
that do not conform to the typical distribution without prior knowledge of the 
outlier model. Our work focuses on nonparametric tests that rely on 
distributional distances rather than parametric assumptions. In particular, we 
analyze kernel-based approaches using the Maximum Mean Discrepancy (MMD)
and information-theoretic measures like the Kullback--Leibler (KL) divergence) 
for distinguishing anomalous data. Theoretical guarantees are provided in terms of error exponents, consistency, and asymptotic optimality.  
\\
\\
To validate the theoretical results, we conduct simulations on synthetic data, with Gaussian distributions as typical benchmarks, examining detection accuracy, 
error probabilities, and stopping criteria. Our study highlights the trade-offs between threshold selection, sample complexity, and computational efficiency in sequential and fixed-sample settings. The results show that kernel methods and 
divergence measures provide reliable tools for anomaly detection in both structured and unstructured data.  

\vspace{\baselineskip}

\noindent
\textbf{Keywords:} Anomaly Detection, Hypothesis Testing, Universal Outlier Problem, Maximum Mean Discrepancy, KL Divergence, Error Exponent, Nonparametric Methods, Sequential Testing, Gaussian Distribution, Stopping Time, Threshold